\subsection{Die Lagrangedichte der Quantenelektrodynamik}
\label{subsec:lagrangedichte_quantenelektrodynamik}

Im vorangegangenen Abschnitt wurde gezeigt, wie die lokale
\(U(1)\)-Eichsymmetrie zur kovarianten Ableitung und damit zur
elektromagnetischen Kopplung führt. Nun sollen das Dirac-Feld, das
elektromagnetische Feld und ihre Wechselwirkung nicht mehr getrennt
nebeneinanderstehen.

Die zentrale Leitfrage dieses Abschnitts lautet daher:

\begin{quote}
	\textbf{Wie lässt sich die gesamte Quantenelektrodynamik -- freies
	Elektronenfeld, freies elektromagnetisches Feld und ihre Wechselwirkung
	-- in einer einzigen mathematischen Struktur zusammenfassen?}
\end{quote}

Diese Aufgabe übernimmt die
Lagrangedichte\index{Lagrangedichte}\index{Quantenelektrodynamik!Lagrangedichte}.
Sie ist der kompakte mathematische Bauplan der Theorie.

\subsubsection*{Ein mathematischer Bauplan für die Theorie}

Eine Feldtheorie muss festlegen, wie sich ihre Felder ohne äußeren
Eingriff entwickeln und wie sie einander beeinflussen. Die
Lagrangedichte \(\mathcal L\) fasst genau diese Informationen für jeden
Raumzeitpunkt zusammen.

Ihre Grundidee lässt sich symbolisch schreiben als

\begin{equation*}
	\text{Lagrangedichte}
	\longrightarrow
	\text{Bewegungsgleichungen und Wechselwirkungen}.
\end{equation*}

Die Lagrangedichte selbst ist keine direkt gemessene Größe. Aus ihr
lassen sich mithilfe des Wirkungsprinzips und der
Euler-Lagrange-Gleichungen für Felder\index{Euler-Lagrange-Gleichungen!Felder}
die Bewegungsgleichungen der Theorie gewinnen. Für den physikalischen
Gedankengang genügt diese Aussage; eine vollständige Variationsrechnung
ist im Haupttext nicht erforderlich.

Im Folgenden werden weiterhin SI-Einheiten verwendet. Damit bleiben die
Faktoren \(c\), \(\hbar\) und \(\mu_0\) sichtbar und die Formeln schließen
unmittelbar an die Konventionen aus
Abschnitt~\ref{subsec:eichsymmetrie_elektromagnetische_wechselwirkung}
an. Natürliche Einheiten würden dieselbe Physik beschreiben, hier aber
lediglich diese Konstanten unterdrücken.

\subsubsection*{Das Dirac-adjungierte Feld}

Um das Dirac-Feld in einer relativistisch geeigneten Lagrangedichte zu
verwenden, benötigt man neben \(\psi\) das
Dirac-adjungierte Feld\index{Dirac-adjungiertes Feld}

\begin{equation}
	\bar{\psi}
	=
	\psi^\dagger\gamma^0.
	\label{eq:dirac_adjunktion_qed}
\end{equation}

Dabei bezeichnet \(\psi^\dagger\) das komplex konjugierte und
transponierte Dirac-Feld. Der zusätzliche Faktor \(\gamma^0\) sorgt
dafür, dass sich Kombinationen wie \(\bar{\psi}\psi\) unter
Lorentz-Transformationen in der benötigten Weise verhalten.
\(\bar{\psi}\) ist also kein neues physikalisches Feld, sondern eine an
die relativistische Spinorstruktur angepasste Schreibweise.

Wie bereits in Abschnitt~\ref{subsec:freie_wechselwirkende_quantenfelder}
erläutert, wird die mathematische Struktur zunächst mit \(\psi\) und
\(\bar{\psi}\) notiert. In der quantisierten Theorie werden die
entsprechenden Größen als Feldoperatoren verstanden.

\subsubsection*{Drei physikalische Beiträge}

Die in Abschnitt~\ref{subsec:freie_wechselwirkende_quantenfelder}
vorbereitete Gliederung wird durch die Lagrangedichte mathematisch
konkret:

\begin{equation}
	\mathcal L_{\mathrm{QED}}
	=
	\mathcal L_{\mathrm D}
	+
	\mathcal L_{\mathrm{EM}}
	+
	\mathcal L_{\mathrm{int}}.
	\label{eq:zerlegung_lagrangedichte_qed}
\end{equation}

Jeder Term erfüllt eine eigene physikalische Aufgabe: Der erste
beschreibt das freie Dirac-Feld, der zweite das freie elektromagnetische
Feld und der dritte die Kopplung zwischen beiden.

\subsubsection*{Das freie Dirac-Feld}

Der erste Beitrag muss die relativistische Bewegung, die Masse und die
Spinstruktur des Elektronenfeldes enthalten. Diese Informationen werden
in der freien Dirac-Lagrangedichte zusammengefasst:

\begin{equation}
	\mathcal L_{\mathrm D}
	=
	\bar{\psi}
	\left(
	\mathrm{i}\hbar c\,\gamma^\mu\partial_\mu
	-
	m_{\mathrm e}c^2
	\right)
	\psi.
	\label{eq:lagrangedichte_freies_diracfeld_qed}
\end{equation}

Der Ableitungsterm beschreibt die Bewegung des Feldes durch die
Raumzeit und enthält über die Gamma-Matrizen seine relativistische
Spinstruktur. Der Massenterm legt die Ruhemasse des Elektrons fest. Aus
diesem Teil der Lagrangedichte folgt die freie Dirac-Gleichung, die
bereits in Abschnitt~\ref{subsec:freie_wechselwirkende_quantenfelder}
eingeführt wurde; sie muss hier nicht erneut hergeleitet werden.

\subsubsection*{Das freie elektromagnetische Feld}

Auch das elektromagnetische Feld benötigt einen eigenen dynamischen
Term. Er muss elektrische und magnetische Felder beschreiben können,
ohne bereits eine Ladung oder einen Strom als Quelle vorauszusetzen.
Mit dem in Abschnitt~\ref{subsec:eichsymmetrie_elektromagnetische_wechselwirkung}
eingeführten Feldstärketensor lautet dieser Beitrag

\begin{equation}
	\mathcal L_{\mathrm{EM}}
	=
	-
	\frac{1}{4\mu_0}
	F_{\mu\nu}F^{\mu\nu}.
	\label{eq:lagrangedichte_freies_em_feld_qed}
\end{equation}

Die Kontraktion \(F_{\mu\nu}F^{\mu\nu}\) verbindet alle Komponenten des
elektrischen und magnetischen Feldes zu einer Lorentz-invarianten
Größe. Der Faktor \(\mu_0\) und das Vorzeichen gehören zur hier
verwendeten SI-Konvention mit der Minkowski-Signatur \((+,-,-,-)\).
Physikalisch erzeugt dieser Term die Dynamik des freien
elektromagnetischen Feldes und damit die Maxwell-Gleichungen ohne
Quellen.

\subsubsection*{Die Kopplung von Materie und elektromagnetischem Feld}

Die beiden freien Terme beschreiben noch keine gegenseitige
Beeinflussung. Entscheidend ist deshalb die Frage, an welcher Stelle
das elektromagnetische Feld in die Dynamik des geladenen Dirac-Feldes
eintritt.

Wie in Abschnitt~\ref{subsec:eichsymmetrie_elektromagnetische_wechselwirkung}
gezeigt wurde, ersetzt die minimale Kopplung im Dirac-Term
\(\partial_\mu\) durch \(D_\mu\). Setzt man
Gleichung~\eqref{eq:kovariante_ableitung_qed} ein, wird ohne lange
Umformung sichtbar:

\begin{equation}
	\begin{aligned}
	&\bar{\psi}
	\left(
	\mathrm{i}\hbar c\,\gamma^\mu D_\mu
	-
	m_{\mathrm e}c^2
	\right)
	\psi
	\\[4pt]
	&\qquad=
	\mathcal L_{\mathrm D}
	-
	qc\,\bar{\psi}\gamma^\mu\psi A_\mu.
	\end{aligned}
	\label{eq:minimale_kopplung_lagrangedichte_qed}
\end{equation}

Die minimale Kopplung verändert also nicht nur ein Symbol in der
Ableitung. Sie verbindet das Dirac-Feld unmittelbar mit dem
Viererpotential und erzeugt dadurch den fundamentalen
Wechselwirkungsterm der QED.

Die Kombination des Dirac-Feldes, die dabei auftritt, besitzt eine
anschauliche Bedeutung. Der elektrische
Viererstrom\index{Viererstrom!Dirac-Feld} ist

\begin{equation}
	j^\mu
	=
	qc\,\bar{\psi}\gamma^\mu\psi
	=
	\left(c\rho,\mathbf j\right).
	\label{eq:viererstrom_diracfeld_qed}
\end{equation}

Die zeitliche Komponente enthält die elektrische Ladungsdichte
\(\rho\), die drei räumlichen Komponenten bilden die elektrische
Stromdichte \(\mathbf j\). Das Vorzeichen der Ladung steckt bereits in
\(q\); für das Elektronenfeld gilt \(q=-e\).

Damit kann der Wechselwirkungsterm besonders übersichtlich geschrieben
werden als

\begin{equation}
	\mathcal L_{\mathrm{int}}
	=
	-j^\mu A_\mu
	=
	-qc\,\bar{\psi}\gamma^\mu\psi A_\mu.
	\label{eq:wechselwirkungsterm_qed}
\end{equation}

Wegen \(A_\mu=(\phi/c,-\mathbf A)\) verbindet die Kontraktion
\(j^\mu A_\mu\) sowohl die Ladungsdichte mit dem elektrischen Potential
als auch die Stromdichte mit dem Vektorpotential:

\begin{equation*}
	j^\mu A_\mu
	=
	\rho\phi
	-
	\mathbf j\cdot\mathbf A.
\end{equation*}

Ein einziger Kopplungsterm beschreibt damit die elementare
Wechselwirkung zwischen dem geladenen Dirac-Feld und dem
elektromagnetischen Feld. Er wurde nicht unabhängig hinzugefügt,
sondern entsteht aus dem in
Abschnitt~\ref{subsec:eichsymmetrie_elektromagnetische_wechselwirkung}
begründeten Übergang
\(\partial_\mu\longrightarrow D_\mu\).

\subsubsection*{Die vollständige QED-Lagrangedichte}

Nun können alle Bestandteile in einer einzigen Formel zusammengeführt
werden:

\begin{equation}
	\begin{aligned}
	\mathcal L_{\mathrm{QED}}
	&=
	\bar{\psi}
	\left(
	\mathrm{i}\hbar c\,\gamma^\mu D_\mu
	-
	m_{\mathrm e}c^2
	\right)
	\psi
	-
	\frac{1}{4\mu_0}F_{\mu\nu}F^{\mu\nu}
	\\[4pt]
	&=
	\mathcal L_{\mathrm D}
	+
	\mathcal L_{\mathrm{EM}}
	+
	\mathcal L_{\mathrm{int}}.
	\end{aligned}
	\label{eq:lagrangedichte_qed}
\end{equation}

In der ersten Zeile steckt der Wechselwirkungsterm bereits in
\(D_\mu\). Die zweite Zeile macht die physikalische Zerlegung sichtbar.
Die Gleichung enthält somit die freie Bewegung des Elektronenfeldes,
die Ausbreitung des elektromagnetischen Feldes und die Kopplung beider
Felder.

\begin{PhysicsBox}[Die gesamte QED in drei Termen]
	\label{box:qed_drei_terme}

	Die Lagrangedichte der Quantenelektrodynamik besitzt die grundlegende
	Struktur

	\begin{equation*}
		\boxed{
		\text{Materiefeld}
		+
		\text{Eichfeld}
		+
		\text{Kopplung}
		}.
	\end{equation*}

	\(\mathcal L_{\mathrm D}\) beschreibt das Elektronenfeld,
	\(\mathcal L_{\mathrm{EM}}\) das elektromagnetische Feld und
	\(\mathcal L_{\mathrm{int}}\) ihre Wechselwirkung. Wenige Terme legen
	damit die physikalische Struktur der gesamten Theorie fest.
\end{PhysicsBox}

\subsubsection*{Von der Lagrangedichte zu physikalischen Vorhersagen}

Wendet man die Euler-Lagrange-Gleichungen auf
\(\mathcal L_{\mathrm{QED}}\) an, erhält man zwei bereits bekannte Seiten
derselben Wechselwirkung:

\begin{itemize}
	\item Die Variation nach dem Dirac-adjungierten Feld \(\bar{\psi}\)
	      liefert die Dirac-Gleichung mit elektromagnetischer Kopplung.
	\item Die Variation nach \(\psi\) führt zur adjungierten Dirac-Gleichung.
	\item Die Variation von \(A_\mu\) liefert die Maxwell-Gleichungen mit
	      \(j^\mu\) als Quelle.
\end{itemize}

Dabei werden \(\psi\) und \(\bar{\psi}\) in der Variationsrechnung als
unabhängige Feldvariablen behandelt.

Das elektromagnetische Feld beeinflusst also die Bewegung des
Dirac-Feldes, während Ladung und Strom des Dirac-Feldes zugleich das
elektromagnetische Feld beeinflussen. Gerade diese wechselseitige
Verknüpfung macht aus zwei freien Feldtheorien die Quantenelektrodynamik.

Die Lagrangedichte ist dennoch nicht selbst das Ergebnis eines
Experiments. Sie ist der Ausgangspunkt, aus dem später
Streuamplituden\index{Streuamplitude},
Wechselwirkungswahrscheinlichkeiten sowie die Emission und Absorption
von Photonen berechnet werden. Ebenso führt sie zu Vorhersagen für
Elektron-Positron-Prozesse und für Präzisionskorrekturen wie die kleine
Abweichung des magnetischen Moments des Elektrons vom Dirac-Wert.

Damit ist das mathematische Fundament gelegt. Für eine konkrete
Berechnung muss nun derjenige Teil der Lagrangedichte genauer betrachtet
werden, der die Felder unmittelbar miteinander verbindet. Die nächste
Leitfrage lautet deshalb:

\begin{quote}
	\textbf{Welcher Teil der QED-Lagrangedichte beschreibt das elementare
	Wechselwirkungsereignis zwischen Elektron, Positron und Photon?}
\end{quote}

Diese Frage führt zum QED-Vertex, dem elementaren Baustein der
Wechselwirkung:

\begin{equation*}
	\mathcal L_{\mathrm{QED}}
	\longrightarrow
	\mathcal L_{\mathrm{int}}
	\longrightarrow
	\text{QED-Vertex}.
\end{equation*}
